% Source: Weinberg GR, physical PDF pp. 372--376
%         (printed pp. 350--354).
% Coverage: Chapter 11 bibliography and references 1--39.
% Figures/tables/footnotes: none.
% Status: source-reviewed and compile-clean.
% Uncertainties: none.

\chapterbackmatter{Bibliography}

\subsubsection*{Relativistic Astrophysics in General}

\begin{itemize}
  \item[\(\square\)] \emph{Quasars and High Energy Astronomy}, ed.\ by
  K.~N.~Douglas, I.~Robinson, A.~Schild, E.~L.~Sch\"ucking,
  J.~A.~Wheeler, and N.~J.~Woolf (Second ``Texas'' Symposium on
  Relativistic Astrophysics, Gordon and Breach, New York, 1969).

  \item[\(\square\)] \emph{High Energy Astrophysics}, ed.\ by L.~Gratton
  (Proceedings of the International School of Physics ``Enrico Fermi,''
  Course XXXV, Academic Press, New York, 1966).

  \item[\(\square\)] \emph{Quasi-Stellar Sources and Gravitational
  Collapse}, ed.\ by I.~Robinson, A.~Schild, and E.~L.~Sch\"ucking
  (First ``Texas'' Symposium on Relativistic Astrophysics, University of
  Chicago Press, Chicago, 1965).

  \item[\(\square\)] Ya.~B.~Zeldovich and I.~D.~Novikov,
  ``Relativistic Astrophysics I,'' \emph{Usp. Fiz. Nauk} \textbf{84}, 377
  (1964) [trans.\ \emph{Soviet Physics Uspekhi}, May--June 1965].

  \item[\(\square\)] Ya.~B.~Zeldovich and I.~D.~Novikov,
  ``Relativistic Astrophysics II,'' \emph{Usp. Fiz. Nauk} \textbf{86}, 447
  (1965) [trans.\ \emph{Soviet Physics Uspekhi} \textbf{8}, 522 (1965)].

  \item[\(\square\)] Ya.~B.~Zeldovich and I.~D.~Novikov,
  \emph{Relativistic Astrophysics: Volume 1, Stars and Relativity},
  translated by K.~S.~Thorne and W.~D.~Arnett (University of Chicago
  Press, Chicago, 1971). I regret that this comprehensive work was not yet
  available during the preparation of this chapter.
\end{itemize}

\subsubsection*{Nonrelativistic Theory of Stellar Structure}

\begin{itemize}
  \item[\(\square\)] S.~Chandrasekhar, \emph{An Introduction to the Study
  of Stellar Structure} (Dover Publications, New York, 1939).

  \item[\(\square\)] E.~E.~Salpeter, ``Stellar Structure Leading up to
  White Dwarfs and Neutron Stars,'' in \emph{Relativity Theory and
  Astrophysics. 3. Stellar Structure}, ed.\ by J.~Ehlers (American
  Mathematical Society, Providence, R.~I., 1967), p.~1.

  \item[\(\square\)] M.~Schwarzschild, \emph{Structure and Evolution of
  the Stars} (Princeton University Press, Princeton, N.~J., 1958).
\end{itemize}

\subsubsection*{Pulsars and Neutron Stars}

\begingroup
\setlength{\emergencystretch}{1em}
\begin{itemize}
  \item[\(\square\)] A.~G.~W.~Cameron, ``Neutron Stars,'' in
  \emph{Annual Review of Astronomy and Astrophysics}, Vol.~8, ed.\ by
  L.~Goldberg (Annual Reviews, Inc., Palo Alto, 1970), p.~179.

  \item[\(\square\)] A.~G.~W.~Cameron, ``How Are Neutron Stars Formed?''
  \emph{Comments Astrophys. and Space Phys.} \textbf{1}, 172 (1969).

  \item[\(\square\)] S.~Frautschi, J.~N.~Bahcall, G.~Steigman, and
  J.~C.~Wheeler, ``Ultradense Matter,'' \emph{Comments Astrophys. and
  Space Phys.}, to be published.

  \item[\(\square\)] V.~L.~Ginzburg, ``Superfluidity and
  Superconductivity in Astrophysics,'' \emph{Comments Astrophys. and Space
  Phys.} \textbf{1}, 81 (1969).

  \item[\(\square\)] T.~Gold, ``The Nature of Pulsars,'' in
  \emph{Contemporary Physics---Trieste Symposium 1968}, ed.\ by A.~Salam,
  Vol.~1 (International Atomic Energy Agency, Vienna, 1969), p.~477.

  \item[\(\square\)] A.~Hewish, ``Pulsars,'' in \emph{Annual Review of
  Astronomy and Astrophysics}, Vol.~8, ed.\ by L.~Goldberg (Annual
  Reviews, Inc., Palo Alto, Cal., 1970), p.~265.

  \item[\(\square\)] L.~D.~Landau and E.~M.~Lifshitz,
  \emph{Statistical Physics}, trans.\ by E.~Peierls and R.~F.~Peierls
  (Pergamon Press, London, 1958), Chapter XI.

  \item[\(\square\)] J.~P.~Ostriker, ``The Nature of Pulsars,''
  \emph{Scientific American}, January 1971, p.~48.

  \item[\(\square\)] M.~A.~Ruderman, ``Solid Stars,''
  \emph{Scientific American}, March 1971, p.~24.

  \item[\(\square\)] Symposium on the Crab Pulsar,
  \emph{Pub. Astron. Soc. Pac.} \textbf{82}, No.~486 (1970).

  \item[\(\square\)] J.~A.~Wheeler, ``Superdense Stars,''
  \emph{Annual Review of Astronomy and Astrophysics}, Vol.~4, ed.\ by
  L.~Goldberg (Annual Reviews, Inc., Palo Alto, Cal., 1966), p.~393.
\end{itemize}
\endgroup

\subsubsection*{Supermassive Objects}

\begin{itemize}
  \item[\(\square\)] R.~V.~Wagoner, ``Physics of Massive Objects,'' in
  \emph{Annual Review of Astronomy and Astrophysics}, Vol.~7, ed.\ by
  L.~Goldberg (Annual Reviews, Inc., Palo Alto, Cal., 1969), p.~553.
\end{itemize}

\subsubsection*{Gravitational Collapse}

\begin{itemize}
  \item[\(\square\)] B.~K.~Harrison, K.~S.~Thorne, M.~Wakano, and
  J.~A.~Wheeler, \emph{Gravitational Theory and Gravitational Collapse}
  (University of Chicago Press, Chicago, 1965).

  \item[\(\square\)] S.~W.~Hawking and D.~W.~Sciama, ``Singularities in
  Collapsing Stars and Expanding Universes,'' \emph{Comments Astrophys.
  and Space Phys.} \textbf{1}, 1 (1969).

  \item[\(\square\)] R.~Geroch, ``Singularities,'' in
  \emph{Relativity---Proceedings of the Relativity Conference in the
  Midwest}, ed.\ by M.~Carmeli, S.~I.~Fickler, and L.~Witten (Plenum
  Press, New York, 1970), p.~259.

  \item[\(\square\)] M.~M.~May and R.~H.~White, ``Hydrodynamic
  Calculations of General Relativistic Collapse,'' in \emph{Relativity
  Theory and Astrophysics. 3. Stellar Structure}, op.\ cit., p.~96.

  \item[\(\square\)] C.~W.~Misner, ``Gravitational Collapse,'' in
  \emph{Astrophysics and General Relativity} (1968 Brandeis University
  Summer Institute in Theoretical Physics), Vol.~1, ed.\ by M.~Chretien,
  S.~Deser, and J.~Goldstein (Gordon and Breach Science Publishers, New
  York, 1969).

  \item[\(\square\)] R.~Penrose, ``On Gravitational Collapse,'' in
  \emph{Contemporary Physics---Trieste Symposium 1968}, ed.\ by A.~Salam,
  Vol.~1 (International Atomic Energy Agency, Vienna, 1969), p.~545.

  \item[\(\square\)] R.~Penrose, ``Structure of Space-Time,'' in
  \emph{Batelle Rencontres}, ed.\ by C.~M.~DeWitt and J.~A.~Wheeler
  (W.~A.~Benjamin, New York, 1968), p.~121.

  \item[\(\square\)] K.~S.~Thorne, ``Nonspherical Gravitational Collapse:
  Does it Produce Black Holes?'' \emph{Comments Astrophys. and Space
  Phys.} \textbf{2}, 191 (1970).

  \item[\(\square\)] R.~Ruffini and J.~A.~Wheeler, ``Introducing the Black
  Hole,'' \emph{Physics Today}, January 1971, p.~30.

  \item[\(\square\)] J.~A.~Wheeler, ``Geometrodynamics and the Issue of
  the Final State,'' in \emph{Relativity, Groups, and Topology}, ed.\ by
  C.~DeWitt and B.~DeWitt (Gordon and Breach Science Publishers, New
  York, 1964), p.~317.
\end{itemize}

For material on the quasi-stellar objects, see the bibliography to Chapter~14.

\chapterbackmatter{References}

\begin{enumerate}
  \item \phantomsection\label{ch11-ref-1}
  B.~K.~Harrison, K.~S.~Thorne, M.~Wakano, and J.~A.~Wheeler,
  \emph{Gravitation Theory and Gravitational Collapse} (University of
  Chicago Press, Chicago, 1965), Appendix B; J.~M.~Bardeen, unpublished
  Ph.D.\ thesis, California Institute of Technology, 1965. For the
  extension of this theorem to slowly rotating stars, see J.~B.~Hartle and
  K.~S.~Thorne, \emph{Ap. J.} \textbf{158}, 179 (1969).

  \item \phantomsection\label{ch11-ref-2}
  Harrison, Thorne, Wakano, and Wheeler, op.\ cit., Chapter 3.

  \item \phantomsection\label{ch11-ref-3}
  See, for example, P.~M.~Morse and H.~Feshbach, \emph{Methods of
  Mathematical Physics} (McGraw--Hill, New York, 1953), p.~278.

  \item \phantomsection\label{ch11-ref-4}
  The Lane--Emden functions are extensively discussed by S.~Chandrasekhar,
  \emph{Stellar Structure} (Dover Publications, New York, 1939),
  Chapter IV.

  \item \phantomsection\label{ch11-ref-5}
  Chandrasekhar, op.\ cit., Table 4.

  \item \phantomsection\label{ch11-ref-6}
  A.~Ritter, \emph{Wiedemann Ann.} \textbf{11}, 332 (1880);
  E.~Betti, \emph{Nuovo Cimento} \textbf{7}, 26 (1880).

  \item \phantomsection\label{ch11-ref-7}
  For a detailed discussion of stellar stability, see P.~Ledoux, in
  \emph{Stars and Stellar Structure VIII: Stellar Structure}, ed.\ by
  L.~H.~Aller and D.~B.~McLaughlin (University of Chicago Press, Chicago,
  1965), Chapter 10. Relativistic effects are considered by
  S.~Chandrasekhar, \emph{Ap. J.} \textbf{140}, 417 (1964).

  \item \phantomsection\label{ch11-ref-8}
  S.~Chandrasekhar, \emph{Mon. Not. Roy. Astron. Soc.} \textbf{95}, 207
  (1935). Also see L.~D.~Landau, \emph{Phys. Z. Sowjetunion} \textbf{1},
  285 (1932).

  \item \phantomsection\label{ch11-ref-9}
  Harrison, Thorne, Wakano, and Wheeler, ref.~1, Figure 5, and Chapter 10.

  \item \phantomsection\label{ch11-ref-10}
  J.~R.~Oppenheimer and G.~M.~Volkoff, \emph{Phys. Rev.} \textbf{55}, 374
  (1939). Some earlier references include L.~Landau, ref.~8;
  W.~Baade and F.~Zwicky, \emph{Proc. Nat. Acad. Sci. U.S.} \textbf{20},
  254 (1934); J.~R.~Oppenheimer and R.~Serber, \emph{Phys. Rev.}
  \textbf{54}, 540 (1938); R.~C.~Tolman, \emph{Phys. Rev.} \textbf{55},
  364 (1939).

  \item \phantomsection\label{ch11-ref-11}
  C.~W.~Misner and H.~S.~Zapolsky, \emph{Phys. Rev. Letters} \textbf{12},
  635 (1964).

  \item \phantomsection\label{ch11-ref-12}
  Harrison, Thorne, Wakano, and Wheeler, ref.~1, Appendix A.

  \item \phantomsection\label{ch11-ref-13}
  Y.~C.~Leung and C.~G.~Wang, to be published. Also see C.~G.~Wang,
  W.~K.~Rose, and S.~L.~Schlenker, \emph{Ap. J.} \textbf{160}, L17
  (1970); H.~Lee, Y.~C.~Leung, and C.~G.~Wang, \emph{Ap. J.}
  \textbf{166}, 387 (1971).

  \item \phantomsection\label{ch11-ref-14}
  S.~Tsuruta and A.~G.~W.~Cameron, \emph{Canadian J. Phys.} \textbf{44},
  1895 (1966).

  \item \phantomsection\label{ch11-ref-15}
  A.~G.~W.~Cameron, \emph{Ann. Rev. Astron. and Astrophys.} \textbf{8},
  179 (1970).

  \item \phantomsection\label{ch11-ref-16}
  M.~Ruderman, \emph{Nature} \textbf{223}, 597 (1969).

  \item \phantomsection\label{ch11-ref-17}
  A.~B.~Migdal, \emph{Nucl. Phys.} \textbf{13}, 655 (1959). For other
  references, see Cameron, ref.~15.

  \item \phantomsection\label{ch11-ref-18}
  The effects of high magnetic fields are considered by V.~Canuto and
  H.~Y.~Chiu, \emph{Phys. Rev.} \textbf{173}, 1210, 1220, 1229 (1968).
  For other references, see Cameron, ref.~15.

  \item \phantomsection\label{ch11-ref-19}
  The effects of rotation are considered by J.~B.~Hartle, \emph{Ap. J.}
  \textbf{150}, 1005 (1967); J.~B.~Hartle and K.~S.~Thorne,
  \emph{Ap. J.} \textbf{153}, 807 (1968); \emph{ibid.} \textbf{158}, 719
  (1969).

  \item \phantomsection\label{ch11-ref-20}
  A.~Hewish, S.~J.~Bell, J.~D.~H.~Pilkington, P.~F.~Scott, and
  R.~A.~Collins, \emph{Nature} \textbf{217}, 709 (1968).

  \item \phantomsection\label{ch11-ref-21}
  T.~Gold, \emph{Nature} \textbf{218}, 731 (1968).

  \item \phantomsection\label{ch11-ref-22}
  For a summary, see A.~Hewish, \emph{Ann. Rev. Astron. and Astrophys.}
  \textbf{8}, 265 (1970).

  \item \phantomsection\label{ch11-ref-23}
  F.~Hoyle and W.~A.~Fowler, \emph{Mon. Not. Roy. Astron. Soc.}
  \textbf{125}, 169 (1963); \emph{Nature} \textbf{197}, 533 (1963);
  F.~Hoyle, W.~A.~Fowler, G.~R.~Burbidge, and E.~M.~Burbidge,
  \emph{Ap. J.} \textbf{139}, 909 (1964).

  \item \phantomsection\label{ch11-ref-24}
  This is obtained by comparison with W.~A.~Fowler, in
  \emph{Quasi-Stellar Sources and Gravitational Collapse} (University of
  Chicago Press, Chicago, 1965), p.~56, Eq.~(24).

  \item \phantomsection\label{ch11-ref-25}
  K.~Schwarzschild, \emph{Sitzungsberichte Preuss. Akad. Wiss.}, 424
  (1916).

  \item \phantomsection\label{ch11-ref-26}
  H.~Bondi, \emph{Proc. Roy. Soc. (London)} \textbf{A281}, 39 (1964);
  also, in \emph{Lectures on General Relativity}, ed.\ by S.~Deser and
  K.~W.~Ford (Prentice--Hall, Englewood Cliffs, New Jersey, 1964),
  p.~375.

  \item \phantomsection\label{ch11-ref-27}
  However, see S.~A.~Bludman and M.~Ruderman, \emph{Phys. Rev.}
  \textbf{170}, 1176 (1968); M.~A.~Ruderman, \emph{Phys. Rev.}
  \textbf{172}, 1286 (1968); S.~A.~Bludman and M.~A.~Ruderman,
  \emph{Phys. Rev.} \textbf{D1}, 3243 (1970).

  \item \phantomsection\label{ch11-ref-28}
  G.~S.~Bisnovatyi-Kogan and Ya.~B.~Zeldovich, \emph{Astrofizika}
  \textbf{5}, 223 (1969). For a discussion of stability of relativistic
  gas spheres and clusters of point masses with arbitrarily large central
  redshifts, see G.~S.~Bisnovatyi-Kogan and K.~S.~Thorne, \emph{Ap. J.}
  \textbf{160}, 875 (1970); E.~D.~Fackerell, J.~R.~Ipser, and
  K.~S.~Thorne, \emph{Comments Astrophys. and Space Phys.} \textbf{1},
  140 (1969).

  \item \phantomsection\label{ch11-ref-29}
  F.~Hoyle and W.~A.~Fowler, \emph{Nature} \textbf{213}, 373 (1967);
  also see H.~S.~Zapolsky, \emph{Ap. J.} \textbf{153}, L163 (1968).

  \item \phantomsection\label{ch11-ref-30}
  G.~Birkhoff, \emph{Relativity and Modern Physics} (Harvard University
  Press, Cambridge, Mass., 1923), p.~253. Also see S.~Deser and
  B.~E.~Laurent, \emph{Am. J. Phys.} \textbf{36}, 789 (1968).

  \item \phantomsection\label{ch11-ref-31}
  R.~C.~Tolman, \emph{Proc. Nat. Acad. Sci. U.S.A.} \textbf{20}, 3
  (1934).

  \item \phantomsection\label{ch11-ref-32}
  J.~R.~Oppenheimer and H.~Snyder, \emph{Phys. Rev.} \textbf{56}, 455
  (1939). For further details on the Oppenheimer--Snyder solution and
  other spherically symmetric solutions, see O.~Klein, in \emph{Werner
  Heisenberg und die Physik unserer Zeit} (Vieweg, Braunschweig, Germany,
  1961); F.~Hoyle, W.~A.~Fowler, G.~R.~Burbidge, and E.~M.~Burbidge,
  \emph{Ap. J.} \textbf{139}, 909 (1964); F.~Hoyle and W.~A.~Fowler, in
  \emph{Quasi-Stellar Sources and Gravitational Collapse}, ed.\ by
  I.~Robinson, A.~Schild, and E.~L.~Sch\"ucking (University of Chicago
  Press, Chicago, 1965); C.~W.~Misner and D.~H.~Sharp, \emph{Phys. Rev.}
  \textbf{136}, B571 (1964); G.~C.~McVittie, \emph{Ap. J.} \textbf{140},
  401 (1964); G.~C.~McVittie, \emph{Ann. Inst. Henri Poincar\'e}
  \textbf{6}, No.~1 (1967); M.~M.~May and R.~H.~White, \emph{Phys. Rev.}
  \textbf{141}, 1232 (1966); S.~A.~Colgate and R.~H.~White,
  \emph{Ap. J.} \textbf{143}, 626 (1966). For asymmetric collapse, see
  J.~M.~Cohen, \emph{Phys. Rev.} \textbf{173}, 1258 (1968);
  M.~Fujimoto, \emph{Ap. J.} \textbf{152}, 523 (1968);
  V.~de la Cruz, J.~E.~Chase, and W.~Israel, \emph{Phys. Rev. Letters}
  \textbf{24}, 423 (1970), and so on.

  \item \phantomsection\label{ch11-ref-33}
  B.~Carter, \emph{Phys. Rev. Letters} \textbf{26}, 331 (1971). For other
  rigorous theorems on the possible forms of exterior metrics under various
  conditions, see A.~Lichnerowicz, \emph{Th\'eories relativistes de la
  gravitation} (Masson, Paris, 1955); S.~Deser, \emph{C. R. Acad. Sci.
  Paris} \textbf{264}, 805 (1967); W.~Israel, \emph{Phys. Rev.}
  \textbf{164}, 1776 (1967); A.~G.~Doroshkevich, Ya.~B.~Zeldovich, and
  I.~D.~Novikov, \emph{Zh. Eksp. Teor. Fiz.} \textbf{49}, 170 (1965)
  [trans.\ \emph{Sov. Phys. JETP} \textbf{22}, 122 (1966)];
  R.~M.~Wald, \emph{Phys. Rev. Letters} \textbf{26}, 1653 (1971), and so
  on.

  \item \phantomsection\label{ch11-ref-34}
  F.~Hoyle and W.~Fowler, \emph{Nature} \textbf{197}, 533 (1963).

  \item \phantomsection\label{ch11-ref-35}
  F.~J.~Dyson, \emph{Comments Astrophys. and Space Phys.} \textbf{1}, 75
  (1969); C.~Leibovitz and W.~Israel, \emph{Phys. Rev.} \textbf{D1}, 3226
  (1970).

  \item \phantomsection\label{ch11-ref-36}
  R.~Penrose, \emph{Riv. Nuovo Cimento} \textbf{1}, Numero Speciale, 252
  (1969).

  \item \phantomsection\label{ch11-ref-37}
  R.~Penrose, \emph{Phys. Rev. Letters} \textbf{14}, 57 (1965);
  S.~W.~Hawking, \emph{Proc. Roy. Soc.} \textbf{294A}, 511 (1966);
  \emph{ibid.} \textbf{295A}, 490 (1966);
  \emph{ibid.} \textbf{300A}, 187 (1967);
  \emph{ibid.} \textbf{308A}, 433 (1967);
  S.~W.~Hawking and R.~Penrose, \emph{Proc. Roy. Soc.} \textbf{314A}, 529
  (1970); and ref.~36.

  \item \phantomsection\label{ch11-ref-38}
  Gravitational radiation from vibrating black holes has been considered by
  W.~H.~Press, \emph{Ap. J.} \textbf{170}, L105 (1971). Gravitational
  radiation from matter falling into black holes has been considered by
  M.~Davis, R.~Ruffini, W.~H.~Press, and R.~K.~Price,
  \emph{Phys. Rev. Letters} \textbf{27}, 1466 (1971).

  \item \phantomsection\label{ch11-ref-39}
  See A.~G.~W.~Cameron, \emph{Nature} \textbf{229}, 178 (1971);
  R.~E.~Wilson, \emph{Ap. J.} \textbf{170}, 529 (1971).
\end{enumerate}
